% ==============================================================================
% PRESENTACIÓN — ITERACIÓN 1 (MDA) · Proyecto AUTOVIDA
% CLENCH Software Development · 16:9
% ==============================================================================
\documentclass[11pt]{article}
\usepackage[utf8]{inputenc}
\usepackage[T1]{fontenc}
\usepackage[spanish,es-tabla,es-nodecimaldot]{babel}
\usepackage[papersize={254mm,143mm}, left=14mm, right=14mm, top=12mm, bottom=11mm]{geometry}
\usepackage{graphicx}
\usepackage[table]{xcolor}
\usepackage{tabularx}
\usepackage{array}
\usepackage{enumitem}
\usepackage{microtype}
\usepackage{tikz}
\usepackage{textcomp}
\usepackage[sfdefault,lf]{FiraSans}
\renewcommand{\familydefault}{\sfdefault}
\usepackage{hyperref}

% ------------------------------------------------------------------ Paleta
\definecolor{tinta}{HTML}{0E2A33}
\definecolor{cian}{HTML}{00A6C7}
\definecolor{cianoscuro}{HTML}{005D75}
\definecolor{magenta}{HTML}{AD0090}
\definecolor{gris}{HTML}{5A6872}
\definecolor{grisclaro}{HTML}{9AA7AE}
\definecolor{banda}{HTML}{EDF5F7}
\definecolor{cebra}{HTML}{F7FAFB}
\definecolor{linea}{HTML}{C9D9DF}

\hypersetup{colorlinks=true, urlcolor=cian, hidelinks}

\pagestyle{empty}
\setlength{\parindent}{0pt}
\linespread{1.1}

\setlength{\arrayrulewidth}{0.7pt}
\arrayrulecolor{cianoscuro}
\renewcommand{\arraystretch}{1.5}
\setlength{\tabcolsep}{6pt}

\newcolumntype{L}{>{\raggedright\arraybackslash}X}
\newcolumntype{D}{>{\raggedleft\arraybackslash}p{2.1cm}}
\newcommand{\cabtxt}[1]{\textbf{\color{tinta}#1}}

% Título de diapositiva: regla corta cian encima, sin eyebrows ni numeración.
\newcommand{\titulo}[1]{%
  {\color{cian}\rule{2.2cm}{2.2pt}}\\[0.30cm]
  {\fontsize{21}{24}\selectfont\bfseries\color{tinta} #1}\\[0.45cm]}

\newcommand{\tituloclaro}[1]{%
  {\color{cian}\rule{2.2cm}{2.2pt}}\\[0.30cm]
  {\fontsize{21}{24}\selectfont\bfseries\color{white} #1}\\[0.45cm]}

% Pie de diapositiva, discreto.
\newcommand{\pie}[1]{%
  \vfill
  {\color{linea}\rule{\textwidth}{0.4pt}}\\[0.12cm]
  {\footnotesize\color{grisclaro} AUTOVIDA\, ·\, Iteración 1\, ·\, CLENCH Software Development
   \hfill #1}}

\newcommand{\pieclaro}[1]{%
  \vfill
  {\color{cianoscuro}\rule{\textwidth}{0.4pt}}\\[0.12cm]
  {\footnotesize\color{gris} AUTOVIDA\, ·\, Iteración 1\, ·\, CLENCH Software Development
   \hfill #1}}

% Viñeta propia: cuadrado cian pequeño.
\setlist[itemize]{leftmargin=1.1em, itemsep=5pt, topsep=2pt,
  label={\raisebox{0.15ex}{\color{cian}\rule{0.42em}{0.42em}}}}

\begin{document}

% ============================================================ 1. PORTADA
\pagecolor{tinta}
\vspace*{0.4cm}
\includegraphics[width=5.6cm]{images/logo_clench.png}

\vspace{1.5cm}

{\fontsize{36}{40}\selectfont\bfseries\color{white} Iteración 1}\\[0.45cm]
{\fontsize{19}{22}\selectfont\color{cian} Proyecto AUTOVIDA}\\[0.3cm]
{\large\color{grisclaro} Agenda digital accesible para alumnado con NEAE}

\vspace{1.1cm}

{\color{grisclaro} 26 de octubre --- 11 de noviembre de 2026\quad·\quad Sprint 1 cerrado}

\pieclaro{Metodologías de Desarrollo Ágil}
\newpage

% ============================================================ 2. OBJETIVO
\pagecolor{white}
\titulo{Qué íbamos a construir}
\vfill

\begin{minipage}[t]{0.56\textwidth}
\vspace{0pt}
\begin{itemize}
  \item Desplegar la \textbf{arquitectura base}: cliente Flutter, API FastAPI y PostgreSQL.
  \item Completar la \textbf{gestión de cuentas, roles y asociaciones} entre alumnado, tutores y personal.
  \item Una primera versión de \textbf{creación y asignación de tareas}.
  \item Dejar operativas las \textbf{vistas de administración}.
\end{itemize}

\vspace{0.5cm}
{\color{gris} No se entrega producto terminado: es un prototipo desechable, pensado para enseñar el planteamiento y abrir la conversación con el centro.}
\end{minipage}%
\hfill
\begin{minipage}[t]{0.38\textwidth}
\vspace{0pt}
\colorbox{banda}{\begin{minipage}{\dimexpr\textwidth-2\fboxsep}
\vspace{2mm}
\hspace{3mm}\begin{minipage}{\dimexpr\textwidth-6mm}
{\footnotesize\color{gris} Alcance previsto}\\[0.25cm]
{\fontsize{30}{32}\selectfont\bfseries\color{tinta} 67}\ {\large\color{gris} PH}\\[0.35cm]
{\footnotesize\color{gris} 38 puntos de las HU-01 a HU-15\\ 29 puntos de despliegue}
\end{minipage}
\vspace{3mm}
\end{minipage}}
\end{minipage}

\pie{2}
\newpage

% ============================================================ 3. LA CIFRA
\pagecolor{tinta}
\tituloclaro{Lo que de verdad pasó}
\vfill

\vspace{0.3cm}

\begin{minipage}[b]{0.46\textwidth}
{\large\color{grisclaro} El alcance creció un}\\[0.2cm]
{\fontsize{76}{78}\selectfont\bfseries\color{magenta} 73\,\%}\\[0.4cm]
{\large\color{white} durante el sprint, y aun así se cerró entero.}
\end{minipage}%
\hfill
\begin{minipage}[b]{0.48\textwidth}
{\small
\arrayrulecolor{cianoscuro}
\begin{tabularx}{\textwidth}{@{}L>{\raggedleft\arraybackslash}p{2.6cm}@{}}
\hline
{\color{white}\textbf{Alcance del sprint}} & {\color{white}\textbf{Puntos}}\\
\hline
{\color{grisclaro}Previsto al planificar} & {\color{white}67}\\
{\color{grisclaro}4 historias adelantadas} & {\color{white}+14}\\
{\color{grisclaro}14 historias nuevas} & {\color{white}+35}\\
\hline
{\color{white}\textbf{Total ejecutado}} & {\color{magenta}\textbf{116}}\\
\hline
\end{tabularx}}
\end{minipage}

\pieclaro{3}
\newpage

% ============================================================ 4. POR QUÉ
\pagecolor{white}
\titulo{Por qué creció: apareció el aula}
\vfill

Durante el sprint se detectó que el modelo no contemplaba \textbf{la clase como entidad}. Sin ella no había forma de que un alumno se buscara dentro de su aula al iniciar sesión, tal como se planteó en clase con el docente.

\vspace{0.55cm}

\begin{minipage}[t]{0.47\textwidth}
\vspace{0pt}
{\color{cianoscuro}\textbf{14 historias nuevas}}\\[0.25cm]
{\small
\begin{itemize}[itemsep=3pt]
  \item Crear, modificar y eliminar clases.
  \item Asignar profesorado y alumnado a un aula.
  \item Crear, consultar, modificar y eliminar tareas de una clase entera.
  \item Cambiar el avatar de un alumno.
\end{itemize}}
\end{minipage}%
\hfill
\begin{minipage}[t]{0.47\textwidth}
\vspace{0pt}
{\color{cianoscuro}\textbf{4 historias adelantadas}}\\[0.25cm]
{\small
\begin{itemize}[itemsep=3pt]
  \item Consultar tareas por fecha: alumno, tutor y personal.
  \item Crear tarea por parte del alumno.
\end{itemize}

\vspace{0.3cm}
{\color{gris} Se adelantaron porque otras historias dependían de ellas para poder avanzar.}}
\end{minipage}

\pie{4}
\newpage

% ============================================================ 5. ENTREGADO
\pagecolor{white}
\titulo{Qué se entregó}
\vfill

\begin{minipage}[t]{0.52\textwidth}
\vspace{0pt}
\begin{itemize}
  \item \textbf{32 historias de usuario} cubiertas por la aplicación, verificadas una a una contra la API en ejecución.
  \item \textbf{Pruebas de aceptación superadas} en backend y en frontend.
  \item Arquitectura completa funcionando en local con un único guion de arranque.
  \item Primera auditoría de accesibilidad del prototipo.
\end{itemize}

\vspace{0.4cm}
{\color{gris} Pendiente: la revisión con el centro. Sus observaciones pueden generar correcciones que se priorizarán para el sprint 2.}
\end{minipage}%
\hfill
\begin{minipage}[t]{0.42\textwidth}
\vspace{0pt}
{\small
\rowcolors{2}{white}{cebra}
\begin{tabularx}{\textwidth}{@{}L>{\raggedleft\arraybackslash}p{1.5cm}@{}}
\hline
\cabtxt{Área} & \cabtxt{Historias}\\
\hline
Administración & 13\\
Personal del centro & 11\\
Tutores y familias & 3\\
Alumnado & 3\\
Autenticación & 2\\
\hline
\textbf{Total} & \textbf{32}\\
\hline
\end{tabularx}}
\end{minipage}

\pie{5}
\newpage

% ============================================================ 6. CÓMO
\pagecolor{white}
\titulo{Cómo nos organizamos}
\vfill

\begin{minipage}[t]{0.46\textwidth}
\vspace{0pt}
{\color{cianoscuro}\textbf{Dos subequipos}}\\[0.2cm]
{\small
\textbf{Frontend} ---\, Pablo Hernández, José Rodríguez y Julián Carrión, que además ejerce de gestor de accesibilidad.\\[0.25cm]
\textbf{Backend} ---\, Yeray Rodríguez, Pablo de la Torre, Miguel Ángel Luque y Amine Azzammouri.}

\vspace{0.45cm}
{\color{cianoscuro}\textbf{Una plantilla común}}\\[0.2cm]
{\small Cada historia se descompone siempre en las mismas \textbf{8 subtareas}: base de datos, endpoint, acceso a datos, diseño de interfaz, implementación, conexión con la API y las dos pruebas de aceptación.}
\end{minipage}%
\hfill
\begin{minipage}[t]{0.48\textwidth}
\vspace{0pt}
{\small
\rowcolors{2}{white}{cebra}
\begin{tabularx}{\textwidth}{@{}L>{\raggedleft\arraybackslash}p{1.6cm}>{\raggedleft\arraybackslash}p{1.5cm}@{}}
\hline
\cabtxt{Integrante} & \cabtxt{Tareas} & \cabtxt{Al día}\\
\hline
José Rodríguez Fernández & 43 & 3,31\\
Julián Carrión Tovar & 43 & 3,31\\
Pablo Hernández Ibáñez & 42 & 3,23\\
Yeray Rodríguez Navas & 32 & 2,46\\
Pablo de la Torre Roldán & 32 & 2,46\\
Miguel Ángel Luque Gómez & 32 & 2,46\\
Amine Azzammouri & 32 & 2,46\\
\hline
\textbf{Total} & \textbf{256} & \textbf{2,81}\\
\hline
\end{tabularx}}

\vspace{0.2cm}
{\footnotesize\color{gris} Ritmo de aceptación de tareas en 13 días laborables, no velocidad en puntos.}
\end{minipage}

\pie{6}
\newpage

% ============================================================ 7. QUÉ FALLÓ
\pagecolor{white}
\titulo{Qué nos costó}
\vfill

\begin{itemize}
  \item \textbf{Integrar al final.} Backend y frontend fusionaron sus ramas tarde, y los conflictos obligaron a convocar reuniones urgentes.
  \item \textbf{Estimar con prisa.} Las 14 historias nuevas se estimaron sobre la marcha, mientras se creaban, sin margen para contrastarlas.
  \item \textbf{Las reuniones híbridas.} Cuando parte del equipo estaba presente y parte conectada, la conversación presencial se imponía y costaba intervenir desde casa.
  \item \textbf{Cuadrar horarios.} Otros proyectos y compromisos personales dificultaron encontrar huecos comunes.
\end{itemize}

\pie{7}
\newpage

% ============================================================ 8. MEJORAS
\pagecolor{white}
\titulo{Qué vamos a cambiar}
\vfill

\begin{minipage}[t]{0.47\textwidth}
\vspace{0pt}
{\color{magenta}\textbf{Integrar de forma continua}}\\[0.25cm]
{\small Registrar en Jira las dependencias y responsables de cada historia \textit{antes} de empezarla, y organizar sesiones de trabajo conjunto ---preferiblemente presenciales--- para fusionar backend y frontend de forma gradual, en lugar de esperar al cierre del sprint.}
\end{minipage}%
\hfill
\begin{minipage}[t]{0.47\textwidth}
\vspace{0pt}
{\color{magenta}\textbf{Reaccionar a los cambios de alcance}}\\[0.25cm]
{\small Cuando se incorpore o adelante una historia, revisar su estimación y medir el impacto en la capacidad del equipo. Se propone un informe que recoja qué se añadió, por qué, cuánto se estimó y qué se decidió.}
\end{minipage}

\vspace{0.7cm}
{\color{gris} Son propuestas, no acuerdos: se validan en la reunión de planificación del sprint 2, donde cada acción recibirá responsable y fecha de revisión.}

\pie{8}
\newpage

% ============================================================ 9. CIERRE
\pagecolor{tinta}
\tituloclaro{Hacia el sprint 2}
\vfill

\vspace{0.4cm}

\begin{minipage}[b]{0.44\textwidth}
{\large\color{grisclaro} Alcance inicial planificado}\\[0.3cm]
{\fontsize{60}{62}\selectfont\bfseries\color{white} 53}\ {\Large\color{grisclaro} PH}\\[0.4cm]
{\color{grisclaro} Frente a los 67 del sprint 1. Un arranque más prudente, a la espera de la sesión de planificación.}
\end{minipage}%
\hfill
\begin{minipage}[b]{0.50\textwidth}
{\color{white} En qué se centra:}\\[0.3cm]
{\small\color{grisclaro}
\begin{itemize}[label={\raisebox{0.15ex}{\color{magenta}\rule{0.42em}{0.42em}}}, itemsep=6pt]
  \item Agenda completa y pasos guiados de actividad.
  \item Ajustes de accesibilidad por perfil de alumno.
  \item Validación con las personas usuarias del centro.
\end{itemize}}
\end{minipage}

\pieclaro{9}

\end{document}
